\documentclass[10pt,a4paper]{amsart}
\usepackage[margin=19mm]{geometry}
\usepackage{amsmath,amssymb,mathtools}
\usepackage[scr=rsfs,cal=euler]{mathalfa}
\usepackage{xfrac}
\usepackage[unicode,hidelinks,bookmarks=false]{hyperref}
\newcommand{\ie}{i.e.\ }
\newcommand{\cf}{cf.\ }
\newcommand{\resp}{respectively\ }
\newcommand{\Subsection}{\subsection}
\providecommand{\nobreakdash}{\nobreak\hbox{-}}
\setlength{\emergencystretch}{2em}
\multlinegap0pt
\usepackage{amssymb}
\usepackage{enumerate}
\usepackage{tikz}
\usepackage{mathtools}
\usepackage{xcolor}
\newtheorem{theorem}{Theorem}
\newtheorem{lemma}{Lemma}
\newcommand{\BV}{\textup{BV}}
\newcommand{\Leb}{\mathcal{L}}
\DeclareMathOperator{\Lip}{\textup{Lip}}
\newcommand{\Lup}{\textup{L}}
\DeclareMathOperator{\Per}{\textup{Per}}
\newcommand{\R}{\mathbb R}
\newcommand{\dx}{\,{\mathrm dx}}
\title{Weighted Sobolev Inequalities via the Meyers--Ziemer Framework: Measures, Isoperimetric Inequalities, and Endpoint Estimates}
\author{Simon Bortz, Kabe Moen, Andrea Olivo, Carlos P\'erez,, Ezequiel Rela}
\date{Source: arXiv:2603.05382v1 (CC BY 4.0)}
\begin{document}
\maketitle

\noindent\emph{Source and licence.} Weighted Sobolev Inequalities via the Meyers--Ziemer Framework: Measures, Isoperimetric Inequalities, and Endpoint Estimates. Version arXiv:2603.05382v1 (CC BY 4.0), distributed by arXiv under the Creative Commons Attribution 4.0 International licence, http://creativecommons.org/licenses/by/4.0/, as recorded in the arXiv licence metadata for this exact version and shown on the abstract page https://arxiv.org/abs/2603.05382v1. Full licence notice: \texttt{licenses/arxiv-2603.05382-CC-BY-4.0.txt}.\par\smallskip

\noindent\textit{Editorial scope.} Counted unit: the article's theorem Gradient (source0.tex lines 532--541). The article's complete original proof of it is delivered with it (source0.tex lines 1427--1483, one \texttt{proof} environment of 57 lines, which the article places away from the statement). No mathematics is written, repaired or re-derived: the statement and the proof are the article's own lines, in the article's own notation, and no retained line is substituted or re-flowed.  Retained uncounted as the source-local context the proof needs: lines 1390--1400; lines 1403--1409.  Nothing was omitted from any retained span, so the package carries no omission record.  No retained line contains a citation, so the wrapper carries no bibliography and no bibliography entry.  No retained line contains a figure environment, an image-inclusion command or drawing code, so no image is delivered and the package declares no asset dependency.  Editorial formatting only, in addition: an AMS \texttt{amsart} preamble that restates the article's own declarations and macros used by the retained lines, namely the environments \texttt{theorem}, \texttt{lemma} on separate counters, together with the standard packages those lines need.  The retained environments print in this document as Theorem 1, Lemma 1 in order of appearance, whereas the article prints the same items with the numbering of its own full document.  The complete native source of the article as distributed in the arXiv e-print for this exact version is bundled unchanged as \texttt{sources/195797/source0.tex}.
\begin{lemma}\label{coarea}
Suppose $w$ is a weight and $u\in \BV(w)$ is non-negative. For $t>0$, set $E_t=\{u>t\}$. Then $E_t$ has finite $w$ perimeter for $\Leb$-a.e. $t\in \R$ and
\[
|Du|_w(\R^n)
=\int_0^\infty \Per(E_t,w)\,\mathrm dt
=\int_0^\infty |\partial E_t|_w(\R^n)\,\mathrm dt
=\int_0^\infty \int_{\R^n} w\,\mathrm d|D\mathbf 1_{E_t}|\,\mathrm dt.
\]
In particular, if $u\in \Lip_c(\R^n)$ then
$$\|\nabla u\|_{\Lup^1(w)}=\int_0^\infty\Per(E_t,w)\,\mathrm dt=\int_0^\infty \int_{\R^n} w\,\mathrm d|D\mathbf 1_{E_t}|\,\mathrm dt.$$
\end{lemma}

\begin{lemma}\label{lowerbdiso}
Suppose $E\subseteq \R^n$ is a bounded set with finite perimeter. Then for any ball $B_r(x)\subseteq \R^n$,
\begin{equation}\label{iso}
\min\{\mathcal L(B_r(x)\cap E),\mathcal L(B_r(x)\cap E^c)\}^{\frac{n-1}{n}}
\le C\,|D\mathbf 1_E|(B_r(x)).
\end{equation}
\end{lemma}

\begin{theorem}\label{Gradient}
Let $\mu$ be a locally finite Borel measure on $\R^n$. Then there exists a dimensional constant $C=C_n$, independent of $\mu$, such that
\begin{equation}\label{newMZthm}
\int_{\R^n} |u| \, \mathrm{d}\mu
\leq
C \int_{\R^n} |\nabla u(x)| \, \textup{M}_1 \mu(x)\, \mathrm{d}x,
\qquad
u \in \Lip_c(\R^n).
\end{equation}
\end{theorem}

\begin{proof}[Proof of Theorem~\ref{Gradient}]
Let $E\subseteq \R^n$ be a bounded open set of finite perimeter. Since $E$ is open, for every $x\in E$ we have
\[
\lim_{r\to\infty}\frac{\mathcal L(B_r(x)\cap E)}{\mathcal L(B_r(x))}=0,
\qquad
\lim_{r\to 0^+}\frac{\mathcal L(B_r(x)\cap E)}{\mathcal L(B_r(x))}=1.
\]
Consequently, there exists $r_x>0$ such that
\[
\frac{\mathcal L(B_{r_x}(x)\cap E)}{\mathcal L(B_{r_x}(x))}=\frac12,
\qquad\text{that is,}\qquad
\mathcal L(B_{r_x}(x)\cap E)=\frac12\,\mathcal L(B_{r_x}(x)).
\]
It follows that
\[
\mathcal L(B_{r_x}(x)\cap E^c)
=\mathcal L(B_{r_x}(x))-\mathcal L(B_{r_x}(x)\cap E)
=\frac12\,\mathcal L(B_{r_x}(x)).
\]
Applying Lemma~\ref{lowerbdiso}, we obtain for each $x\in E$ the estimate
\[
r_x^{n-1}\le C\,|D\mathbf 1_E|(B_{r_x}(x)).
\]

Now let $u\in \Lip_c(\R^n)$. For each $t>0$, the level set $E_t=\{|u|>t\}$ is open, bounded, and has finite perimeter. Let $K\subseteq E_t$ be compact. Then
\[
K\subseteq \bigcup_{x\in K} B_{r_x}(x).
\]
By compactness and the Vitali covering lemma, there exists a finite disjoint collection of balls $B_j=B_{r_j}(x_j)$, with $x_j\in K$, such that
\[
K\subseteq \bigcup_{j=1}^N 3B_j.
\]
Therefore,
\begin{align*}
\mu(E_t)
&\le \sum_{j=1}^N \mu(3B_j)
= c_n\sum_{j=1}^N \frac{\mu(3B_j)}{r_j^{\,n-1}}\, r_j^{\,n-1} \\
&\le c_n\sum_{j=1}^N \frac{\mu(3B_j)}{r_j^{\,n-1}}\,
|D\mathbf 1_{E_t}|(B_j)
\le c_n\sum_{j=1}^N \int_{B_j} \textup{M}_1\mu\,\mathrm d|D\mathbf 1_{E_t}| \\
&\le c_n\int_{\R^n} \textup{M}_1\mu\,\mathrm d|D\mathbf 1_{E_t}|
= c_n\,|\partial E_t|_{\textup{M}_1\mu}(\R^n),
\end{align*}
where we used that for every $x\in B_j$,
\[
\textup{M}_1\mu(x)\ge c_n\,\frac{\mu(3B_j)}{r_j^{\,n-1}}.
\]

Finally, integrating over $t\in(0,\infty)$ and applying the coarea formula with weight $w=\textup{M}_1\mu$ (Lemma~\eqref{coarea}), we obtain
\[
\int_{\R^n} |u|\,\mathrm d\mu
= \int_0^\infty \mu(E_t)\,\mathrm dt
\le C\int_0^\infty \int_{\R^n} \textup{M}_1\mu\,\mathrm d|D\mathbf 1_{E_t}|\,\mathrm dt
\le C\int_{\R^n} |\nabla u|\,\textup{M}_1\mu\,\dx.
\]
This concludes the proof.
\end{proof}

\end{document}
